\documentclass{article}
\usepackage{amsmath,amssymb,amsthm}
\usepackage[hidelinks]{hyperref}
\usepackage{cleveref}

\newtheorem{theorem}{Theorem}[section]
\newtheorem{lemma}[theorem]{Lemma}
\newtheorem{corollary}[theorem]{Corollary}
\theoremstyle{definition}
\newtheorem{definition}[theorem]{Definition}
\numberwithin{equation}{section}

\title{A Note on Sums of Two Squares}
\author{A. Author\\ Department of Mathematics}
\date{}

\begin{document}
\maketitle

\begin{abstract}
We give a short, self-contained account of which primes are sums of two
squares. The argument uses only the pigeonhole principle and elementary
congruences.
\end{abstract}

\section{Introduction}

Our goal is \cref{thm:main}. The key input is \cref{lem:thue}, and the
statement for general integers is \cref{cor:general}.

\begin{theorem}\label{thm:main}
An odd prime $p$ is a sum of two squares if and only if $p \equiv 1
\pmod 4$.
\end{theorem}

\section{Preliminaries}

\begin{definition}\label{def:rep}
An integer $n$ is \emph{represented} if $n = x^2 + y^2$ for some integers
$x$ and $y$.
\end{definition}

Squares are congruent to $0$ or $1$ modulo $4$, so
\begin{equation}\label{eq:mod4}
  x^2 + y^2 \not\equiv 3 \pmod 4
\end{equation}
for all integers $x, y$.

\begin{lemma}[Thue]\label{lem:thue}
Let $p$ be prime and $a$ an integer not divisible by $p$. There are
integers $x, y$ with $0 < |x|, |y| < \sqrt{p}$ and
$ax \equiv y \pmod p$.
\end{lemma}

\begin{proof}
Let $m = \lfloor \sqrt{p} \rfloor$. The $(m+1)^2 > p$ numbers $au - v$ with
$0 \le u, v \le m$ cannot be distinct modulo $p$, and subtracting two equal
ones gives the required $x$ and $y$.
\end{proof}

\section{Proof of the main theorem}

\begin{proof}[Proof of \cref{thm:main}]
Necessity is \eqref{eq:mod4}. Conversely, if $p \equiv 1 \pmod 4$ then
$-1$ is a square modulo $p$, say $a^2 \equiv -1$. Applying
\cref{lem:thue} to $a$ gives
\begin{equation}\label{eq:thue}
  x^2 + y^2 \equiv x^2 + a^2x^2 \equiv 0 \pmod p,
\end{equation}
with $0 < x^2 + y^2 < 2p$. Hence $x^2 + y^2 = p$.
\end{proof}

\begin{corollary}\label{cor:general}
A positive integer is represented in the sense of \cref{def:rep} if and
only if every prime $q \equiv 3 \pmod 4$ divides it to an even power.
\end{corollary}

\begin{proof}
Combine \cref{thm:main} with the identity
$(a^2+b^2)(c^2+d^2) = (ac-bd)^2 + (ad+bc)^2$ and with \eqref{eq:thue}.
\end{proof}

\end{document}
